\subsection{A stationary circulation in the polymer reaction graph}
\label{sec:polymer-circulation}

There is also a nonreversible extension on the original unlabelled
polymer state space.  It uses only the existing merge and split edges.
The modification is state dependent; it is not a kernel depending
only on the sizes of the reacting clusters.
Antisymmetric perturbations preserving equilibrium are an established
mechanism for accelerating the asymptotic convergence rate; see, for
example, Hwang, Hwang-Ma and Sheu~\cite{HwangHwangMaSheu2005} for the
diffusion setting.  The construction below verifies that mechanism on
an explicit cycle of the present polymer reaction graph.

\begin{lemma}\label{lem:simple-polymer-gap}
For every $n\ge2$ and $\beta>0$, the eigenvalue $\gamma_M=\gamma_Q$
of $-L_M$ is simple.  Its real eigenfunction can be written $g(K(m))$,
with $g(k+1)-g(k)>0$ for $1\le k<n$ after a choice of sign.
\end{lemma}

\begin{proof}
The positive ground state $h$ of the count-derivative operator is
the discrete derivative of a first nonconstant count eigenfunction,
so the latter is strictly monotone.  That count eigenvalue is simple.
We check that no higher sector can tie it.

The minimum-ratio argument in Lemma~\ref{lem:ground-ratio} gives the strict inequality
$h(i+1)/h(i)>1/2$.
Indeed, if a minimum ratio $r\le1/2$ occurred at $i$, the two adjacent
eigenvalue equations would give
\[
 0\ge\beta(1-r)+\frac{2i}{n}(r^{-1}-1)-\frac2n>0.
\]
At $r=1/2$ the last expression is
$\beta/2+2(i-1)/n>0$; smaller $r$ only increases it.
The endpoint conventions are the same as in that proposition.
In its higher-sector comparison the extra term
$(\ell-1)\beta[2h(r+1)/h(r)-1]$ is consequently strictly positive
for every retained degree $\ell\ge2$.
Thus all these sectors have bottom eigenvalue strictly above
$\gamma_Q$.  A polymer observable has no degree-one component.
Equality in its Poincar\'e inequality therefore forces its lift
to lie in the one-dimensional first nonconstant count eigenspace.
For $n=2,3$ there are no retained higher sectors to check.
\end{proof}

For $n\ge4$, consider the four distinct polymer configurations
\[
 x=(2,1^{n-2}),\qquad
 u=(2,2,1^{n-4}),\qquad
 v=(4,1^{n-4}),\qquad
 w=(3,1^{n-3}),
\]
written as multisets of cluster sizes.
Every edge of the directed cycle $x\to u\to v\to w\to x$ is an
allowed reaction.  Let $J(a,b)$ be $1$ on its forward edges,
$-1$ on its reverse edges, and $0$ elsewhere.
For an oriented edge $a,b$, write
$C(a,b)=\pi_M(a)q(a,b)=\pi_M(b)q(b,a)$.
Choose
\begin{equation}\label{eq:circulation-size}
 0<|\varepsilon|<
       \min\{C(x,u),C(u,v),C(v,w),C(w,x)\}
\end{equation}
and define
\begin{equation}\label{eq:circulation-rates}
 q_\varepsilon(a,b)
   =q(a,b)+\varepsilon\frac{J(a,b)}{\pi_M(a)}
 \quad(a\ne b).
\end{equation}
Since $\sum_bJ(a,b)=0$, the diagonal generator entries are unchanged.

\begin{proposition}\label{prop:polymer-circulation}
The rates \eqref{eq:circulation-rates} define an irreducible
nonreversible polymer chain with invariant law $\pi_M$.
Its symmetric part is exactly $L_M$.
Its asymptotic spectral rate is strictly greater than $\gamma_M$,
whereas its best prefactor-one $L^2(\pi_M)$ contraction exponent
is exactly $\gamma_M$.
\end{proposition}

\begin{proof}
Condition \eqref{eq:circulation-size} keeps every original reaction
rate positive.  The antisymmetry and zero divergence of $J$ show that
$\pi_M$ is stationary.  The stationary forward and backward fluxes
on each cycle edge differ by $2\varepsilon$, so detailed balance fails.
Write $L_\varepsilon=L_M+\varepsilon A$, where
$A(a,b)=J(a,b)/\pi_M(a)$ for $a\ne b$ and $A(a,a)=0$.
Then $A1=0$ and $A^*=-A$, proving the symmetric-part claim.

Let $g(K)$ be the eigenfunction in Lemma~\ref{lem:simple-polymer-gap},
normalized in $L^2(\pi_M)$.  The two states $u,w$ have cardinality
$n-2$, while $x,v$ have cardinalities $n-1,n-3$.
Consequently
\[
 (Ag)(u)=
   \frac{g(n-3)-g(n-1)}{\pi_M(u)}\ne0.
\]
Also $\langle g,Ag\rangle_{\pi_M}=0$ by skew-adjointness.

For a nonzero eigenvalue $\lambda$ of $-L_\varepsilon$ with complex
eigenvector $F$, stationarity implies $\mathbb E_{\pi_M}F=0$.
Therefore
\[
 \operatorname{Re}\lambda
 =\frac{\langle F,-L_MF\rangle_{\pi_M}}{\|F\|_2^2}
 \ge\gamma_M.
\]
Equality would put $F$ in the complex span of $g$, by simplicity of
the bottom eigenspace.  But
$-L_\varepsilon g=\gamma_Mg-\varepsilon Ag$ is not in that span.
Equality is impossible.  Since the spectrum is finite, its least
nonzero real part is strictly larger than $\gamma_M$.

The symmetric part remains $L_M$, so differentiation of the squared
$L^2$ norm gives contraction at exponent $\gamma_M$.
Differentiating at $t=0$ with $F=g$ rules out any larger exponent
with prefactor one.
\end{proof}

The small-circulation improvement has a positive explicit
second-order coefficient.  Put $B=-L_M$, let $g$ have norm one, and
let $R$ be the inverse of $B-\gamma_M I$ on
$\{1,g\}^{\perp}$.
The lowest nonzero eigenvalue branch satisfies
\begin{equation}\label{eq:circulation-perturbation}
 \lambda(\varepsilon)
 =\gamma_M+\varepsilon^2\langle Ag,R\,Ag\rangle_{\pi_M}
     +O(\varepsilon^4),
 \qquad \langle Ag,R\,Ag\rangle_{\pi_M}>0.
\end{equation}
Indeed, simple-eigenvalue perturbation gives first correction
$R\,Ag$ to the eigenvector and the displayed quadratic coefficient.
The branch is real near zero.  Since
$L_{-\varepsilon}=L_\varepsilon^*$, that real branch is even in
$\varepsilon$, giving the stated remainder.
The coefficient is positive because $Ag\ne0$ and $R$ is positive
definite on $\{1,g\}^{\perp}$.
The separated spectrum at $\varepsilon=0$ ensures that this branch
is the asymptotic rate for sufficiently small $\varepsilon$.

This construction proves strict acceleration for each fixed
$n,\beta$.  It does not provide an improvement bounded away from
zero uniformly in $n$: the admissible stationary current in
\eqref{eq:circulation-size} and the coefficient in
\eqref{eq:circulation-perturbation} depend on the chosen cycle.
